% !TeX program = pdflatex
% !TeX root = ./sprawko.tex
\documentclass[11pt,a4paper]{article}

\usepackage[T1]{fontenc}
\usepackage[utf8]{inputenc}
\usepackage[polish]{babel}
\usepackage{geometry}
\usepackage{graphicx}
\usepackage{tabularx}
\usepackage{array}
\usepackage{caption}
\usepackage{url}
\usepackage{placeins}
\usepackage{pgfplots}
\usepackage{hyperref}

\pgfplotsset{compat=1.18}
\geometry{margin=1.8cm}
\setlength{\parindent}{0pt}
\setlength{\parskip}{0.45em}
\renewcommand{\arraystretch}{1.35}
\captionsetup{font=footnotesize,labelfont=bf}
\hypersetup{colorlinks=true,linkcolor=black,urlcolor=blue}
\setcounter{topnumber}{3}
\setcounter{bottomnumber}{2}
\setcounter{totalnumber}{4}
\renewcommand{\topfraction}{0.92}
\renewcommand{\bottomfraction}{0.82}
\renewcommand{\textfraction}{0.05}
\renewcommand{\floatpagefraction}{0.75}

\newcommand{\pole}[1]{\textit{#1}\par\vspace{0.8cm}}
\newcommand{\wynik}[2]{%
  \begin{figure}[!htbp]
    \centering
    \includegraphics[width=\linewidth,height=0.66\textheight,keepaspectratio]{#1/#2}
    \caption{\protect\path{#2}}
  \end{figure}
}
\newcommand{\notatka}[2]{%
  \clearpage
  \thispagestyle{plain}
  \begin{center}
    \includegraphics[width=\linewidth,height=0.82\textheight,keepaspectratio]{zdj_z_zaj/#1}
    \par\vspace{0.25cm}
    {\small #2}
  \end{center}
}

\begin{document}

\begin{tabularx}{\textwidth}{|>{\raggedright\arraybackslash}p{0.17\textwidth}|X|>{\raggedright\arraybackslash}p{0.19\textwidth}|}
\hline
\pole{Grupa lab.} & \pole{18.11.2026} & \pole{25.11.2026}\\
\hline
\multicolumn{3}{|l|}{\pole{11 - iperf}}\\
\hline
\end{tabularx}

\begin{tabularx}{\textwidth}{|X|>{\raggedright\arraybackslash}p{0.32\textwidth}|}
\pole{Filip Będkowski, Danylo Drygo} & \pole{}\\
\hline
\end{tabularx}

\vspace{0.4cm}
\begin{center}
  {\Large\bfseries Część praktyczna}
\end{center}

\section{Opis wykonanego ćwiczenia}

\subsection{Sprzęt i oprogramowanie}

Do prób lokalnych wykorzystaliśmy dwa komputery z systemem Linux: stację
kliencką (adres \texttt{192.168.102.57}) oraz serwer iperf3
(\texttt{192.168.102.6}). Urządzenia łączyła badana skrętka Ethernet; po
pierwszej serii wymieniliśmy przewód i powtórzyliśmy pomiary. Do prób
zewnętrznych użyliśmy serwerów \texttt{iperf3.moji.fr},
\texttt{ping.online.net} i \texttt{speedtest.milkywan.fr}.

Przepustowość TCP i UDP mierzyliśmy programem \texttt{iperf3}. Dla TCP
wykonywaliśmy testy domyślne oraz z rozmiarami okna \texttt{100K}, \texttt{210K},
\texttt{1M} i \texttt{3M}; weryfikowaliśmy również kierunek odwrotny opcją
\texttt{-R}. Testy trwały 15 s. Odpowiedzi DNS sprawdzaliśmy programem
\texttt{dig}: dla domen agh.edu.pl,
facebook.com, google.com i onet.pl, przez resolver domyślny
(\texttt{192.168.102.1}), Google DNS (\texttt{8.8.8.8}) i Quad9
(\texttt{9.9.9.9}).

\subsection{Schemat ćwiczenia}

\begin{center}
\fbox{\begin{minipage}{0.94\linewidth}
\centering
\textbf{Pomiary lokalne}\\
Klient Ubuntu (\texttt{192.168.102.57})
$\longleftrightarrow$ badana skrętka
$\longleftrightarrow$ serwer iperf3 (\texttt{192.168.102.6})\\[0.35em]
\textbf{Pomiary zewnętrzne}\\
Klient Ubuntu $\longrightarrow$ Internet
$\longrightarrow$ \texttt{iperf3.moji.fr}, \texttt{ping.online.net}
lub \texttt{speedtest.milkywan.fr}\\[0.35em]
\textbf{Pomiary DNS}\\
Klient Ubuntu z \texttt{dig}
$\longrightarrow$ \texttt{192.168.102.1}, \texttt{8.8.8.8}
lub \texttt{9.9.9.9}
\end{minipage}}
\end{center}

\subsection{Wyniki -- zestawienie}

Poniższe przepływności TCP odczytaliśmy z końcowych wierszy \texttt{receiver}
na zrzutach. W tabeli lokalnej zestawiliśmy serię przed wymianą przewodu
(\texttt{test\_serwer\_iperf}) i serię po wymianie (\texttt{test\_nasza\_skretka}).

\begin{table}[!htbp]
\centering
\caption{Lokalne pomiary TCP, przepływność odbiorcza}
\begin{tabularx}{0.92\linewidth}{|X|r|r|}
\hline
\textbf{Ustawienie okna TCP} &
\textbf{Przed wymianą [Mbit/s]} &
\textbf{Po wymianie [Mbit/s]}\\
\hline
Domyślne & 902,0 & 92,9\\
\hline
\texttt{-w 100K} & 811,0 & 92,9\\
\hline
\texttt{-w 210K} & 810,0 & 94,1\\
\hline
\end{tabularx}
\end{table}

\begin{figure}[!htbp]
\centering
\begin{tikzpicture}
\begin{axis}[
  width=0.86\linewidth,
  height=7cm,
  xlabel={Rozmiar okna TCP [KiB]},
  ylabel={Przepływność odbiorcza [Mbit/s]},
  xmin=90, xmax=220,
  ymin=0, ymax=1000,
  xtick={100,210},
  ytick={0,200,400,600,800,1000},
  grid=major,
  legend style={at={(0.5,-0.23)},anchor=north,legend columns=2},
]
\addplot+[mark=*,thick] coordinates {(100,811) (210,810)};
\addlegendentry{Przed wymianą skrętki}
\addplot+[mark=square*,thick] coordinates {(100,92.9) (210,94.1)};
\addlegendentry{Po wymianie skrętki}
\end{axis}
\end{tikzpicture}
\caption{Wpływ jawnie ustawionego okna TCP na lokalną przepływność.}
\end{figure}

\begin{table}[!htbp]
\centering
\caption{Przykładowe domyślne pomiary TCP na serwerach publicznych}
\begin{tabularx}{0.92\linewidth}{|X|r|}
\hline
\textbf{Serwer i port} & \textbf{Odbiorca [Mbit/s]}\\
\hline
\texttt{iperf3.moji.fr:5200} & 258\\
\hline
\texttt{ping.online.net:5200} & 489\\
\hline
\texttt{speedtest.milkywan.fr:9202} & 419\\
\hline
\end{tabularx}
\end{table}

\begin{table}[!htbp]
\centering
\caption{Przykładowy wynik zapytania A dla agh.edu.pl}
\begin{tabularx}{\linewidth}{|X|l|r|r|}
\hline
\textbf{Resolver} & \textbf{Odpowiedź IPv4} & \textbf{TTL [s]} & \textbf{Czas [ms]}\\
\hline
Domyślny (\texttt{192.168.102.1}) & \texttt{149.156.96.150} & 84106 & 0\\
\hline
Google DNS (\texttt{8.8.8.8}) & \texttt{149.156.96.150} & 17349 & 32\\
\hline
Quad9 (\texttt{9.9.9.9}) & \texttt{149.156.96.150} & 31230 & 11\\
\hline
\end{tabularx}
\end{table}

Zrzuty obejmują również zapytania A, AAAA, MX, NS i ANY dla badanych domen,
a także testy UDP. W części publicznej serwer MilkyWan był zajęty na porcie
9200, dlatego połączenie zestawiliśmy na porcie 9202. Wybrane próby z oknem
\texttt{1M} zakończyły się błędem i nie zostały potraktowane jako pomiary
przepływności.

\section{Dokumentacja wyników iperf3}

\subsection{Serwer iperf3.moji.fr}
\wynik{iperf3_moji_fr}{01_iperf3-c-iperf3.moji.fr-p5200-t15.png}
\wynik{iperf3_moji_fr}{02_iperf3-c-iperf3.moji.fr-p5200-t15-w100K.png}
\wynik{iperf3_moji_fr}{03_iperf3-c-iperf3.moji.fr-p5200-t15-w210K.png}
\wynik{iperf3_moji_fr}{04_iperf3-c-iperf3.moji.fr-p5200-t15-w1M-error-seria.png}
\wynik{iperf3_moji_fr}{05_iperf3-c-iperf3.moji.fr-p5200-t15-w100K-powtorzenie.png}
\wynik{iperf3_moji_fr}{06_iperf3-c-iperf3.moji.fr-p5200-t15-R-w100K.png}
\wynik{iperf3_moji_fr}{07_iperf3-c-iperf3.moji.fr-p5200-t15-R-w210K.png}
\wynik{iperf3_moji_fr}{08_iperf3-c-iperf3.moji.fr-p5200-t15-R-w210K-powtorzenie.png}
\wynik{iperf3_moji_fr}{09_iperf3-c-iperf3.moji.fr-p5202-t15-R-u.png}
\FloatBarrier

\subsection{Serwer ping.online.net}
\wynik{ping_online_net}{01_dig-ping.online.net-A__iperf3-c-ping.online.net-p5200-t15.png}
\wynik{ping_online_net}{02_iperf3-c-ping.online.net-p5200-t15-w100K.png}
\wynik{ping_online_net}{03_iperf3-c-ping.online.net-p5200-t15-w210K.png}
\wynik{ping_online_net}{04_iperf3-c-ping.online.net-p5200-t15-w1M.png}
\wynik{ping_online_net}{05_iperf3-c-ping.online.net-p5200-t15-w3M.png}
\wynik{ping_online_net}{06_iperf3-c-ping.online.net-p5200-t15-R.png}
\wynik{ping_online_net}{07_iperf3-c-ping.online.net-p5200-t15-R-w100K.png}
\wynik{ping_online_net}{08_iperf3-c-ping.online.net-p5200-t15-R-w210K.png}
\wynik{ping_online_net}{09_iperf3-c-ping.online.net-p5200-t15-R-w1M.png}
\wynik{ping_online_net}{10_iperf3-c-ping.online.net-p5200-t15-R-w3M.png}
\wynik{ping_online_net}{11_iperf3-c-ping.online.net-p5200-t15-u-R.png}
\wynik{ping_online_net}{12_iperf3-UDP-serwery-ping-ping6-ping90ms-p5200-p5208-t15-proby.png}
\FloatBarrier

\subsection{Serwer speedtest.milkywan.fr}
\wynik{speedtest_milkywan}{01_iperf3-c-speedtest.milkywan.fr-p9200-t15-busy-p9202-t15.png}
\wynik{speedtest_milkywan}{02_iperf3-c-speedtest.milkywan.fr-p9202-t15-w100K.png}
\wynik{speedtest_milkywan}{03_iperf3-c-speedtest.milkywan.fr-p9202-t15-w210K.png}
\wynik{speedtest_milkywan}{04_iperf3-c-speedtest.milkywan.fr-p9202-t15-w1M.png}
\wynik{speedtest_milkywan}{05_iperf3-c-speedtest.milkywan.fr-p9202-t15-w3M.png}
\wynik{speedtest_milkywan}{06_iperf3-c-speedtest.milkywan.fr-p9202-t15-R.png}
\wynik{speedtest_milkywan}{07_iperf3-c-speedtest.milkywan.fr-p9202-t15-R-w100K.png}
\wynik{speedtest_milkywan}{08_iperf3-c-speedtest.milkywan.fr-p9202-t15-R-w210K.png}
\wynik{speedtest_milkywan}{09_iperf3-c-speedtest.milkywan.fr-p9202-t15-R-w1M.png}
\wynik{speedtest_milkywan}{10_iperf3-c-speedtest.milkywan.fr-p9202-t15-R-w3M.png}
\wynik{speedtest_milkywan}{11_iperf3-c-speedtest.milkywan.fr-p9202-t15-u.png}
\wynik{speedtest_milkywan}{12_iperf3-c-speedtest.milkywan.fr-p9202-t15-u-R.png}
\FloatBarrier

\subsection{Połączenie lokalne -- test naszej skrętki}
\wynik{test_nasza_skretka}{01_iperf3-c-192.168.102.6-t15.png}
\wynik{test_nasza_skretka}{02_iperf3-c-192.168.102.6-t15-w100K.png}
\wynik{test_nasza_skretka}{03_iperf3-c-192.168.102.6-t15-w210K.png}
\wynik{test_nasza_skretka}{04_iperf3-c-192.168.102.6-t15-R.png}
\wynik{test_nasza_skretka}{05_iperf3-c-192.168.102.6-t15-R-w100K.png}
\wynik{test_nasza_skretka}{06_iperf3-c-192.168.102.6-t15-R-w210K.png}
\wynik{test_nasza_skretka}{07_iperf3-c-192.168.102.6-t15-R-w1M-error.png}
\wynik{test_nasza_skretka}{08_iperf3-c-192.168.102.6-t15-R-w210K-powtorzenie.png}
\wynik{test_nasza_skretka}{09_iperf3-c-192.168.102.6-t15-R-u.png}
\FloatBarrier

\subsection{Połączenie lokalne -- test serwera iperf}
\wynik{test_serwer_iperf}{01_iperf3-c-192.168.102.6-t15.png}
\wynik{test_serwer_iperf}{02_iperf3-c-192.168.102.6-t15-w100K.png}
\wynik{test_serwer_iperf}{03_iperf3-c-192.168.102.6-t15-w210K.png}
\wynik{test_serwer_iperf}{04_iperf3-c-192.168.102.6-t15-R.png}
\wynik{test_serwer_iperf}{05_iperf3-c-192.168.102.6-t15-R-powtorzenie.png}
\wynik{test_serwer_iperf}{06_iperf3-c-192.168.102.6-t15-R-w100K.png}
\wynik{test_serwer_iperf}{07_iperf3-c-192.168.102.6-t15-R-w210K.png}
\wynik{test_serwer_iperf}{08_iperf3-c-192.168.102.6-t15-R-w210K-powtorzenie.png}
\wynik{test_serwer_iperf}{09_iperf3-c-192.168.102.6-t15-R-u.png}
\FloatBarrier

\section{Wyniki zapytań DNS}

\subsection{Domena agh.edu.pl}
\wynik{zdj_drygo}{01_dig_agh.edu.pl_A-serwer-domyslny.jpeg}
\wynik{zdj_drygo}{02_dig_agh.edu.pl_A-8.8.8.8.jpeg}
\wynik{zdj_drygo}{03_dig_agh.edu.pl_A-9.9.9.9.jpeg}
\wynik{zdj_drygo}{04_dig_agh.edu.pl_AAAA-9.9.9.9.jpeg}
\wynik{zdj_drygo}{05_dig_agh.edu.pl_ANY-serwer-domyslny.jpeg}
\wynik{zdj_drygo}{06_dig_agh.edu.pl_ANY-8.8.8.8.jpeg}
\wynik{zdj_drygo}{07_dig_agh.edu.pl_ANY-9.9.9.9-TCP.jpeg}
\FloatBarrier

\subsection{Domena facebook.com}
\wynik{zdj_drygo}{08_dig_facebook.com_A-serwer-domyslny.jpeg}
\wynik{zdj_drygo}{09_dig_facebook.com_A-8.8.8.8.jpeg}
\wynik{zdj_drygo}{10_dig_facebook.com_A-9.9.9.9.jpeg}
\wynik{zdj_drygo}{11_dig_facebook.com_AAAA-serwer-domyslny.jpeg}
\wynik{zdj_drygo}{12_dig_facebook.com_AAAA-8.8.8.8.jpeg}
\wynik{zdj_drygo}{13_dig_facebook.com_ANY-9.9.9.9.jpeg}
\wynik{zdj_drygo}{14_dig_facebook.com_MX-serwer-domyslny.jpeg}
\wynik{zdj_drygo}{15_dig_facebook.com_MX-8.8.8.8.jpeg}
\wynik{zdj_drygo}{16_dig_facebook.com_NS-serwer-domyslny.jpeg}
\wynik{zdj_drygo}{17_dig_facebook.com_NS-8.8.8.8.jpeg}
\FloatBarrier

\subsection{Domena google.com}
\wynik{zdj_drygo}{18_dig_google.com_A-8.8.8.8.jpeg}
\wynik{zdj_drygo}{19_dig_google.com_A-9.9.9.9.jpeg}
\wynik{zdj_drygo}{20_dig_google.com_ANY-serwer-domyslny-TCP.jpeg}
\wynik{zdj_drygo}{21_dig_google.com_ANY-8.8.8.8.jpeg}
\wynik{zdj_drygo}{22_dig_google.com_ANY-9.9.9.9.jpeg}
\wynik{zdj_drygo}{23_dig_google.com_MX-serwer-domyslny.jpeg}
\wynik{zdj_drygo}{24_dig_google.com_MX-8.8.8.8.jpeg}
\FloatBarrier

\subsection{Domena onet.pl}
\wynik{zdj_drygo}{25_dig_onet.pl_A-serwer-domyslny.jpeg}
\wynik{zdj_drygo}{26_dig_onet.pl_A-9.9.9.9.jpeg}
\wynik{zdj_drygo}{27_dig_onet.pl_AAAA-8.8.8.8.jpeg}
\wynik{zdj_drygo}{28_dig_onet.pl_AAAA-9.9.9.9.jpeg}
\wynik{zdj_drygo}{29_dig_onet.pl_ANY-serwer-domyslny.jpeg}
\wynik{zdj_drygo}{30_dig_onet.pl_ANY-8.8.8.8.jpeg}
\wynik{zdj_drygo}{31_dig_onet.pl_ANY-9.9.9.9.jpeg}
\FloatBarrier

\section{Opracowanie wyników}

\subsection{Wpływ okna TCP i wymiany przewodu}

Najwyższą lokalną przepływność uzyskaliśmy przed wymianą skrętki: w teście
domyślnym odbiornik raportował 902 Mbit/s. Po wymianie przewodu wynik spadł do
92,9 Mbit/s, czyli był około 9,7 raza niższy (spadek o około 89,7\%). Co ważne,
po zmianie przewodu wyniki dla okien 100K i 210K pozostały w przedziale
92,9--94,1 Mbit/s. Taki poziom jest bliski praktycznemu maksimum łącza
100 Mbit/s, dlatego obserwacja jest zgodna z wynegocjowaniem niższej szybkości
Ethernet po zmianie przewodu lub jego podłączenia. Sam pomiar iperf3 nie
potwierdza jednak szybkości negocjacji; do jej jednoznacznego sprawdzenia
potrzebny byłby odczyt stanu interfejsu sieciowego.

Przed wymianą skrętki ustawienie 100K dało 811 Mbit/s, a 210K -- 810 Mbit/s,
wobec 902 Mbit/s w teście domyślnym. Różnica pomiędzy dwoma jawnie ustawionymi
oknami była znikoma, a zwiększenie okna z 100K do 210K nie podniosło wyniku.
Po wymianie przewodu przepustowość zatrzymała się w okolicy 93 Mbit/s niezależnie
od tych dwóch wartości. Wskazuje to, że w tej serii ograniczeniem było przede
wszystkim łącze fizyczne, a nie samo ustawienie okna. Próby z oknem 1M, które
zakończyły się komunikatem o nieprawidłowym rozmiarze bufora gniazda, pominęliśmy
na wykresie, ponieważ nie dostarczyły poprawnego wyniku.

\subsection{Porównanie serwerów publicznych i kierunków}

W testach domyślnych uzyskaliśmy około 258 Mbit/s z iperf3.moji.fr,
489 Mbit/s z ping.online.net i 419 Mbit/s z speedtest.milkywan.fr. Różnica
pomiędzy najniższym i najwyższym wynikiem wynosiła niemal 1,9 raza, mimo że
testy trwały po 15 s. Nie traktujemy jej jako różnicy przepustowości samego
łącza dostępowego: pomiary obejmowały inne trasy, odległości, serwery i warunki
ich obciążenia. W trybie odwrotnym \texttt{-R} przepływność również się
zmieniała, co pokazuje, że kierunek transmisji ma znaczenie i nie należy
utożsamiać wyniku wysyłania z wynikiem pobierania.

Rozmiar okna wpływał wyraźniej w testach zdalnych niż po wymianie skrętki.
Przy małym oknie TCP nadawca może mieć mniej niepotwierdzonych danych w locie;
przy większym opóźnieniu na trasie może to ograniczyć osiąganą przepływność.
Zwiększenie okna nie gwarantuje jednak wzrostu wyniku w każdym pojedynczym
teście, ponieważ równocześnie zmieniają się obciążenie serwera, retransmisje
i warunki sieciowe. Na serwerze MilkyWan połączenia na porcie 9200 wielokrotnie
zwracały komunikat o zajętym serwerze; dlatego do dalszych pomiarów wykorzystaliśmy
port 9202. Takich komunikatów nie interpretowaliśmy jako zerowej przepustowości.

\subsection{TCP, UDP i wiarygodność pomiaru}

W lokalnych testach UDP uzyskaliśmy około 1,05 Mbit/s, co odpowiada domyślnemu
docelowemu strumieniowi UDP programu iperf3. W obu zrzutach lokalnych odebrano
1358 z 1358 datagramów (0\% strat); zarejestrowany jitter wynosił około
0,058 ms przed wymianą i 0,035 ms po wymianie przewodu. Wynik potwierdza
poprawny odbiór datagramów w krótkiej, 15-sekundowej próbie, ale nie jest
pomiarem maksymalnej przepustowości łącza. W części publicznych testów UDP
występowały błędy zajętego serwera lub brak połączenia, więc nie łączyliśmy
tych nieudanych prób z wynikami poprawnych transmisji.

Zwracaliśmy uwagę na różnicę między przepływnością nadawcy i odbiorcy oraz na
retransmisje widoczne w wynikach TCP. Przykładowo, w lokalnej próbie domyślnej
przed wymianą skrętki odbiornik raportował 902 Mbit/s, a nadawca 903 Mbit/s;
równocześnie raport zawierał 5190 retransmisji. Powtórzenia i odczyt obu stron
są zatem ważne przy ocenie stabilności, a pojedynczy wynik nie powinien być
podstawą do wniosku o jakości całego łącza.

\subsection{Odpowiedzi DNS, czasy i TTL}

Dla zapytania A o agh.edu.pl wszystkie trzy resolvery zwróciły ten sam adres
\texttt{149.156.96.150}. W zarejestrowanym przykładzie czas odpowiedzi wynosił
0 ms dla resolvera domyślnego, 32 ms dla Google DNS i 11 ms dla Quad9.
Jednocześnie TTL w odpowiedziach wynosił odpowiednio 84106 s, 17349 s i
31230 s. Zestawienie pokazuje, że zgodność adresu nie oznacza identycznego
stanu pamięci podręcznej: resolvery mogły pobrać rekord w różnym momencie.
Wartość 0 ms dla resolvera lokalnego jest zgodna z szybkim odczytem z cache,
ale pojedynczy pomiar nie pozwala ustalić tego z całkowitą pewnością.

Wyniki kolejnych zapytań DNS nie są stałymi parametrami domeny. Czas odpowiedzi
zależy między innymi od trasy do resolvera i od tego, czy odpowiedź znajduje
się już w jego cache; TTL zmniejsza się podczas przechowywania rekordu i po
odświeżeniu zaczyna od wartości nadanej przez źródło autorytatywne. Odpowiedzi
A i AAAA dotyczą różnych rodzin adresów, rekordy MX wskazują serwery pocztowe,
a NS -- serwery nazw. Zapytanie ANY zwracało zestaw dostępnych odpowiedzi, ale
nie należy traktować go jako gwarantowanego, kompletnego wykazu strefy.
Zmieniające się adresy A dla facebook.com w kolejnych zrzutach są zgodne
z rotacją odpowiedzi lub rozkładaniem ruchu; nie świadczą same w sobie o błędzie
resolvera.

\subsection{Ograniczenia}

W przekazanych materiałach wynikowych udokumentowaliśmy pomiary iperf3 oraz
zapytania dig. Nie ma w nich liczbowych wyników dnsperf, curl, mtr ani
speedtest.pl, dlatego nie dopisywaliśmy dla tych narzędzi rezultatów, których
nie można potwierdzić na podstawie zrzutów. Część publicznych prób zakończyła
się błędem lub komunikatem o zajętym serwerze; zostały zachowane jako
dokumentacja przebiegu, lecz wyłączone z porównań przepływności.

\section*{Wnioski}

\begin{enumerate}
  \item \textbf{Potwierdziliśmy}, że po wymianie skrętki lokalna przepływność
        TCP spadła z około 902 do 92,9 Mbit/s. Stabilne wyniki 92,9--94,1 Mbit/s
        są zgodne z ograniczeniem do łącza 100 Mbit/s; przyczynę należałoby
        dodatkowo zweryfikować stanem negocjacji interfejsu i okablowania.
  \item Dla jawnych okien 100K i 210K przed wymianą uzyskaliśmy odpowiednio
        811 i 810 Mbit/s, a po wymianie 92,9 i 94,1 Mbit/s. W badanym lokalnym
        połączeniu zmiana okna nie kompensowała ograniczenia przepustowości
        powstałego po zmianie przewodu.
  \item Wyniki publicznych serwerów różniły się od 258 do 489 Mbit/s w testach
        domyślnych. Uwzględniliśmy odległość, obciążenie i kierunek transmisji,
        dlatego nie utożsamiamy wyniku pojedynczego serwera z szybkością
        gwarantowaną przez nasze łącze.
  \item W lokalnym UDP uzyskaliśmy około 1,05 Mbit/s, jitter poniżej 0,06 ms
        i brak utraconych datagramów w zarejestrowanych próbach. Do oceny UDP
        potrzebne są jednocześnie przepływność, jitter i procent strat.
  \item Dla agh.edu.pl różne resolvery zwróciły ten sam adres IPv4, ale inne
        wartości TTL i czasu odpowiedzi. Wyniki są zgodne z niezależnym
        buforowaniem rekordów DNS; dlatego porównywaliśmy domenę, typ rekordu
        i resolver, a nie samą wartość TTL.
  \item Powtórzenia i odrzucanie nieudanych prób zwiększają wiarygodność
        wniosków. Wnioski o narzędziach bez zachowanych wyników pomiarów
        pozostawiliśmy poza zakresem sprawozdania.
\end{enumerate}

\section*{Załącznik: notatki z zajęć}

\notatka{02_notatki_dns_ttl_i_rekordy.jpg}{Notatki: system DNS, TTL i typy rekordów.}
\notatka{03_notatki_wyniki_pomiarow_dns.jpg}{Zapis wyników i TTL dla badanych domen.}
\notatka{04_notatki_analiza_ttl_i_odpowiedzi_dns.jpg}{Notatki do analizy odpowiedzi DNS.}
\notatka{05_notatki_scenariusz_serwery_iperf3_i_dns.jpg}{Scenariusz pomiarów publicznych i adresy resolverów.}
\notatka{06_notatki_scenariusz_testy_lokalne_curl_mtr.jpg}{Notatki do testów lokalnych i narzędzi sieciowych.}
\notatka{01_notatki_iperf3_dns_strona-tytulowa.jpg}{Strona z danymi grupy i nazwiskami.}

\end{document}
